% !TeX root = ../../Book1.tex
% 纲要标题：### 15.3 正规扩张
% 纲要位置：计划纲要.md，第 542 行。

\section{正规扩张}
\label{b1:sec:1503}


% 纲要内容（仅作写作计划，尚非教材正文）：
%
% * 分裂条件与嵌入条件；代数扩张到自身的嵌入自动为自同构
% * 正规闭包；有限生成元最小多项式的分裂域与全部嵌入像的生成域
% * 正规性对中间域和塔的行为
%

分裂域包含一个方程的所有根。正规扩张把这一要求推广为：凡是不可约方程已有一个根落入扩张，其余根也必须一并落入。用嵌入来表述时，这意味着改变根的选择不会把整个扩张带到另一个子域。本节的扩张均为代数扩张。

\ProseParagraphBreak
给定代数扩张 \(L/K\)，由\cref{b1:thm:algebraic-closure}可取 \(L\) 的代数闭包 \(\Omega\)。因为 \(\Omega/L\) 与 \(L/K\) 都代数，\cref{b1:thm:algebraic-tower}说明 \(\Omega/K\) 代数；\(\Omega\) 又是代数闭域，所以它同时是包含指定 \(L\) 的 \(K\) 的代数闭包。下文的根与嵌入均在这样选定的 \(\Omega\) 中讨论。

\subsection{分裂条件与嵌入条件}
\label{b1:subsec:1503:01}

\begin{definition}\label{b1:def:normal-extension}
设 \(L/K\) 为代数扩张。如果每个在 \(L\) 中有根的 \(K\) 系数不可约多项式都在 \(L\) 中分裂，则称 \(L/K\) 为\term[zhenggui-kuozhang]{正规扩张}[normal extension]。
\end{definition}
这里的分裂允许重根：只要求多项式成为一次因子的乘积，不要求这些因子互不相同。正规性关心所有根是否留在域内；不同根的个数是下一章可分性所研究的问题。

\begin{proposition}\label{b1:prop:algebraic-self-embedding}
设 \(L/K\) 为代数扩张，\(\sigma:L\rightarrow L\) 为逐点固定 \(K\) 的域嵌入。则 \(\sigma\) 为 \(L\) 的 \(K\) 自同构，无须假设 \(L/K\) 正规。
\end{proposition}
\begin{proof}
即使 \(L/K\) 无限，每个元素仍受一个有限次数的方程约束。给定 \(a\in L\)，令 \(R\) 为 \(m_{a,K}\) 在 \(L\) 中的不同根组成的集合。它有限且包含 \(a\)；\(\sigma\) 固定多项式的系数，故将 \(R\) 单射地映入自身。有限集上的单射必为双射，所以存在 \(b\in R\) 使 \(\sigma(b)=a\)。每个 \(a\in L\) 都有原像，故 \(\sigma\) 满射。
\end{proof}
本命题的目标域已经是 \(L\)，所以只须解决满射性。正规性的嵌入判据还须保证：当目标扩大为代数闭包时，像不会离开 \(L\)。

\begin{theorem}\label{b1:thm:normal-embedding-criterion}
设 \(L/K\) 为代数扩张，\(\Omega\) 为包含 \(L\) 的 \(K\) 的代数闭包。以下条件等价：
\begin{enumerate}
\item \(L/K\) 正规。
\item 每个保持 \(K\) 的嵌入 \(\sigma:L\rightarrow\Omega\) 都满足 \(\sigma(L)=L\)。
\item 存在一族 \(K[x]\) 中的非零多项式，使 \(L\) 在 \(K\) 上由它们在 \(\Omega\) 中的全部根生成。
\end{enumerate}
有限扩张时，这些条件还等价于 \(L\) 是某一个非零 \(K\) 系数多项式的分裂域。
\end{theorem}
\begin{proof}
先设正规。对任意 \(a\in L\)，其最小多项式在 \(L\) 中分裂，\(\sigma(a)\) 也是该多项式的根，所以 \(\sigma(L)\subseteq L\)。将目标限制为 \(L\) 后，\cref{b1:prop:algebraic-self-embedding}给出满射性，故 \(\sigma(L)=L\)。这里适用的是逐元素的有限根集论证，既不要求 \(L/K\) 有限，也没有把一般无限集上的单射当作满射。

设第二个条件成立。若不可约 \(f\in K[x]\) 在 \(L\) 中有根 \(a\)，要证明它分裂，就须把任意另一个根 \(b\in\Omega\) 也纳入 \(L\)。\cref{b1:lem:simple-embedding-extension}给出 \(K(a)\rightarrow K(b)\)、\(a\mapsto b\) 的 \(K\) 同构，再由\cref{b1:thm:embedding-extension}延拓为 \(L\rightarrow\Omega\)。其像按假设为 \(L\)，故 \(b\in L\)。因此全部根属于 \(L\)，证明正规性。

正规时，取 \(L\) 中所有元素在 \(K\) 上的最小多项式；这些非零多项式的全部根都在 \(L\) 中，并在 \(K\) 上生成 \(L\)，得到第三个条件。反过来，若 \(L\) 在 \(K\) 上由一族非零多项式在 \(\Omega\) 中的全部根生成，则每个 \(K\) 嵌入把每个有限根集单射地映入自身，因而置换该根集。生成域遂被映到自身且映满，满足第二个条件。

若 \(L/K\) 有限，取有限生成元 \(a_1,\ldots,a_r\)，令 \(f=\prod_{i=1}^r m_{a_i,K}\)。正规性使各最小多项式都在 \(L\) 中分裂，故全部根生成的域包含于 \(L\)；这个域又包含各 \(a_i\)，因而包含 \(L\)。所以 \(L\) 正是 \(f\) 的分裂域；反方向已由第三个条件得到。
\end{proof}
第三个条件必须取所选多项式的全部根。仅由其中一些根生成的域可以是任意代数扩张，未必正规。例如 \(\mathbb Q(\sqrt2)/\mathbb Q\) 正规；\(\mathbb Q(\sqrt[3]2)/\mathbb Q\) 不正规，因为 \(x^3-2\) 的另两个根不实。此处三次扩张和六次分裂域之间的差别，正是正规性所检测的差别。

\subsection{正规闭包}
\label{b1:subsec:1503:02}

\begin{proposition}\label{b1:prop:normal-closure}
设 \(L/K\) 为代数扩张，\(\Omega\) 为包含 \(L\) 的 \(K\) 的代数闭包。则在所选 \(\Omega\) 中，按子域的包含关系存在包含 \(L\) 的最小正规扩张 \(N/K\)。它由所有 \(a\in L\) 的最小多项式在 \(\Omega\) 中的全部根生成，称为 \(L/K\) 在 \(\Omega\) 中的\term[zhenggui-bibao]{正规闭包}[normal closure]。它也有描述
\[
N=K\left(\bigcup_{\sigma\in\operatorname{Hom}_K(L,\Omega)}\sigma(L)\right),
\]
其中 \(\operatorname{Hom}_K(L,\Omega)\) 表示逐点固定 \(K\) 的嵌入 \(L\rightarrow\Omega\) 全体。若 \(L/K\) 有限且 \(L=K(a_1,\ldots,a_r)\)，则 \(N\) 正是 \(\prod_{i=1}^r m_{a_i,K}\) 在 \(\Omega\) 中的分裂域，特别地 \(N/K\) 有限。
\end{proposition}
\begin{proof}
按所述根集生成 \(N\)，\cref{b1:thm:normal-embedding-criterion}的第三个条件直接保证 \(N/K\) 正规，且每个 \(a\in L\) 都在生成根集中，故 \(L\subseteq N\)。若正规扩张 \(M/K\)、\(L\subseteq M\subseteq\Omega\)，则每个 \(a\) 的最小多项式在 \(M\) 中分裂，全部生成根都属于 \(M\)，故 \(N\subseteq M\)。

记所有像域生成的域为 \(M_0\)。每个 \(K\) 嵌入将 \(a\in L\) 送到 \(m_{a,K}\) 的一个根，故 \(\sigma(L)\subseteq N\)，从而 \(M_0\subseteq N\)。反过来，任取 \(m_{a,K}\) 在 \(\Omega\) 中的根 \(\beta\)。\cref{b1:lem:simple-embedding-extension}给出 \(a\mapsto\beta\) 的 \(K(a)\rightarrow\Omega\) 嵌入，再由\cref{b1:thm:embedding-extension}延拓到 \(L\)。于是 \(\beta\) 属于某个像域，故 \(N\) 的全部生成根都属于 \(M_0\)。这证明 \(N=M_0\)：正规闭包把 \(L\) 在共同代数闭包中的所有 \(K\) 同构像一起纳入。

若 \(L=K(a_1,\ldots,a_r)\) 为有限扩张，取 \(m_{a_i,K}\) 的乘积在 \(\Omega\) 内的分裂域 \(N_0\)。它有限、正规并包含 \(L\)，故 \(N\subseteq N_0\)。另一方面，\(N/K\) 正规且包含各 \(a_i\)，所以包含各 \(m_{a_i,K}\) 的全部根，故 \(N_0\subseteq N\)。因此 \(N=N_0\) 为有限扩张；这个具体描述不依赖生成集的选择。
\end{proof}
例如 \(\mathbb Q(\sqrt[3]2)\) 的正规闭包就是 \(\mathbb Q(\sqrt[3]2,\zeta)\)，其中 \(\zeta^2+\zeta+1=0\)、\(\zeta\neq1\)。这里先固定共同代数闭包，才可以谈“最小子域”；若不固定环境，结论应改写成保持基域和已给扩张的同构意义下唯一。

\subsection{正规性与中间域、扩张塔}
\label{b1:subsec:1503:03}

\begin{proposition}\label{b1:prop:normal-tower-behavior}
设 \(K\subseteq E\subseteq L\) 为代数扩张塔。若 \(L/K\) 正规，则 \(L/E\) 正规。在同一代数闭包内，任意一族正规 \(K\) 扩张的交以及它们生成的合成域都对 \(K\) 正规。但 \(L/K\) 正规不保证 \(E/K\) 正规，\(E/K\) 与 \(L/E\) 都正规也不保证 \(L/K\) 正规。
\end{proposition}
\begin{proof}
每个 \(E\) 嵌入 \(L\rightarrow\Omega\) 也是 \(K\) 嵌入，故由\cref{b1:thm:normal-embedding-criterion}保持 \(L\)，从而 \(L/E\) 正规。

若 \(a\) 属于一族正规域之交，其在 \(K\) 上的最小多项式在每个域中分裂，所以全部根属于交域，证明交域正规。对于合成域，每个被合成的域由某族 \(K\) 系数多项式的根生成；把各族合并即可满足\cref{b1:thm:normal-embedding-criterion}的第三个条件。

第一个反例取 \(L=\mathbb Q(\sqrt[3]2,\zeta)\)、\(E=\mathbb Q(\sqrt[3]2)\)、\(K=\mathbb Q\)，其正规与非正规已在本节说明。第二个取
\[
K=\mathbb Q,\quad E=\mathbb Q(\sqrt2),\quad L=\mathbb Q(\sqrt[4]2).
\]
\(E\) 是 \(x^2-2\) 的分裂域，而 \(L/E\) 是 \(x^2-\sqrt2\) 的分裂域，故两层正规。但 \(x^4-2\) 在 \(\mathbb Q\) 上不可约，其非实根不在实域 \(L\) 中，故总扩张不正规。
\end{proof}
交和合成是在同一环境中组合若干扩张，域塔则是在一个扩张内改变基域；前两种运算保持正规性，并不意味着正规性沿塔传递。将来用子群描述中间域时，这些不同的行为会分别对应不同的群论条件。
